\documentclass[sigconf,screen]{acmart}

\usepackage{amsmath,amsfonts}
\usepackage{algorithmic}
\usepackage{graphicx}
\usepackage{textcomp}
\usepackage{xcolor}
\usepackage{booktabs}
\usepackage{multirow}
\usepackage{balance}

\begin{document}

\title{ConfTest: Confidence-Calibrated Selective Regression Test Selection via Uncertainty-Aware Deep Ensembles}

\author{Antigravity Research Team}
\affiliation{%
  \institution{Department of Computer Science and Engineering, APJ Abdul Kalam Technological University}
  \city{Thiruvananthapuram}
  \country{India}
}
\email{bbipin@ktu.ac.in}

\begin{abstract}
Regression Test Selection (RTS) accelerates Continuous Integration (CI/CD) pipelines by executing only the subset of tests likely affected by incoming code commits. However, existing Machine Learning (ML)-based RTS techniques suffer from poor probability calibration and overconfidence on out-of-distribution refactorings, causing silent regression escapes. In this paper, we propose \textbf{ConfTest}, an uncertainty-aware, confidence-calibrated selective regression test selection framework. ConfTest extracts a 32-dimensional canonical feature vector spanning AST syntactic structure, static call-graph coupling, code churn, and historical failure telemetry. We employ a 5-seed deep ensemble to quantify epistemic uncertainty ($\sigma$) and apply post-hoc Temperature Scaling, which reduces Expected Calibration Error from $0.0390$ to $0.0242$ on a held-out temporal validation split (paired improvement $0.0148$, 95\% CI $[0.0030, 0.0393]$). ConfTest introduces a selective prediction policy that executes budget-matched test subsets when confident and falls back to full-suite execution upon detecting elevated epistemic divergence or high-churn refactorings. Our findings are mixed, and we report them as such. On $183$ held-out commits ($86{,}469$ commit--test pairs, labels obtained by execution) the shipped operating point attains \textbf{$100.0\%$ failure recall with zero escaped commits}, but it abstains on $97.81\%$ of commits and therefore reduces test executions by only $3.11\%$, with a wall-clock reduction of $0.0\%$. Tightening the abstention threshold to reach a $32.85\%$ reduction drops held-out recall to $87.69\%$ with $15$ escaped commits, an $8.33$ percentage-point degradation from the validation split and a failure of our pre-registered $95\%$ recall floor. We conclude that the selective-prediction scaffolding is sound and reusable, while the underlying ranker (PR-AUC $0.1393$ against a $3.71\%$ positive rate) is not yet accurate enough to convert abstention into useful savings---a negative result we believe is more useful to the RTS community than the alternative of tuning on the test split.
\end{abstract}

\keywords{Regression Test Selection, Confidence Calibration, Uncertainty Estimation, Continuous Integration, Machine Learning for Software Engineering}

\maketitle

\section{Introduction}
Modern software development relies heavily on Continuous Integration (CI) and Continuous Deployment (CD). However, as codebases grow, executing full regression test suites on every pull request becomes economically prohibitive, consuming thousands of compute runner minutes and blocking developer feedback loops.

Machine Learning-based Regression Test Selection (ML-RTS) frames test selection as a binary classification problem: predicting whether test $t_j$ will fail on commit $c_i$. While promising, conventional ML-RTS models suffer from three fundamental deficiencies:
\begin{enumerate}
    \item \textbf{Miscalibration:} Tree-based models output scores that do not match empirical failure frequencies. On our validation split the uncalibrated ensemble has expected calibration error $0.0390$ and worst-bin error $0.1235$, so a confidence threshold set from raw scores does not mean what it appears to mean.
    \item \textbf{Epistemic Blindness:} Standard models cannot distinguish between confident predictions and predictions on out-of-distribution (OOD) architectural refactorings.
    \item \textbf{Uncontrolled Regression Escapes:} ML-RTS can skip failing tests unless it includes a conservative fallback, risking production outages.
\end{enumerate}

To resolve these challenges, we present \textbf{ConfTest}, a calibrated selective prediction system that combines ML-based ranking with an empirical full-suite fallback policy. Its safety results are observed on held-out data, not formal guarantees for future commits or unseen repositories.

\section{System Architecture \& Methodology}
\subsection{Canonical 32-Feature Mining Pipeline}
For each commit-test pair $(c_i, t_j)$, ConfTest extracts 32 continuous features across four semantic domains:
\begin{itemize}
    \item \textbf{Diff \& Churn (12):} Lines added/deleted, file counts, commit intent indicators.
    \item \textbf{AST Complexity (6):} Function/class counts, cyclomatic complexity delta, assertions.
    \item \textbf{Static Call-Graph (6):} Direct imports, NetworkX shortest path depth, coupling coefficients.
    \item \textbf{Historical Telemetry (8):} Lifetime failure rates, recent-10 failure rates, flakiness scores.
\end{itemize}

\subsection{Deep Ensemble Epistemic Uncertainty Estimation}
We train a 5-seed Deep Ensemble $\mathcal{E} = \{f_{\theta_1}, \dots, f_{\theta_5}\}$ using LightGBM. The ensemble mean failure probability $\bar{p}$ and epistemic uncertainty $\sigma$ are defined as:
\begin{equation}
    \bar{p}(c_i, t_j) = \frac{1}{M} \sum_{m=1}^M f_{\theta_m}(c_i, t_j)
\end{equation}
\begin{equation}
    \sigma(c_i, t_j) = \sqrt{\frac{1}{M} \sum_{m=1}^M \left( f_{\theta_m}(c_i, t_j) - \bar{p}(c_i, t_j) \right)^2}
\end{equation}

\subsection{Post-Hoc Temperature Scaling Calibration}
To ensure predicted probabilities match empirical failure likelihoods, we optimize a scalar temperature parameter $T > 0$ on the validation split by minimizing Negative Log-Likelihood (NLL):
\begin{equation}
    \hat{p}(c_i, t_j) = \sigma_{\text{sigmoid}}\left( \frac{z(c_i, t_j)}{T} \right)
\end{equation}

We search $\log T$ with a bounded scalar minimiser, which converges in $17$ iterations to
$T^* = 0.7941$ (validation NLL $0.19290 \rightarrow 0.18345$). $T^* < 1$ means our ensemble is
\emph{under}-confident on this split, the opposite of the overconfidence usually reported for deep
networks. We also fitted isotonic regression and rejected it: it attains a lower mean ECE
($0.0166$) but raises worst-bin error from $0.1235$ to $0.2098$ (paired degradation $+0.0863$,
$95\%$ CI $[0.0299, 0.1135]$). An abstention gate is thresholded on the tail of the confidence
distribution, so a method that is better on average and much worse in its worst bin is the wrong
trade.

\subsection{Selective Prediction Policy}
ConfTest operates under a cost-optimal risk-coverage policy. For a commit $c_i$:
\begin{equation}
    \text{Mode}(c_i) = \begin{cases}
        \text{SAFE\_FULL\_SUITE} & \text{if } \sigma(c_i) > \tau_{\text{abstain}} \lor \max \hat{p} < \tau_{\text{conf}} \lor \text{OOD}(c_i) \\
        \text{FAST\_SELECTED} & \text{otherwise}
    \end{cases}
\end{equation}

\section{Empirical Evaluation \& Results}

\subsection{Dataset and Protocol}
Ground truth is obtained by execution rather than by proxy. We generate semantic mutants in five
open-source Python repositories, run each repository's suite against each mutant, and record the
per-test outcome; a (commit, test) pair is positive when that test fails on that mutant and passes
on the unmutated baseline. This yields $561{,}711$ labelled pairs over $1{,}213$ commits.

Splits are chronological and cut on mutant boundaries, so no mutant contributes rows to more than
one split: $849$ commits / $391{,}825$ rows for training ($5.017\%$ positive), $181$ / $83{,}417$
for validation ($5.488\%$), and $183$ / $86{,}469$ held out for test ($3.710\%$). The temperature
and every policy threshold are fitted on validation only; the held-out split is scored once per
reported configuration. Confidence intervals are bootstrap intervals over $2{,}000$ resamples with
the \emph{commit} as the resampling unit, because the $86{,}469$ rows are not independent.
\begin{table}[t]
\centering
\caption{Empirical Benchmark Comparison Against 8 RTS Baselines}
\label{tab:baselines}
\small
\begin{tabular}{lcccc}
\toprule
\textbf{RTS Strategy} & \textbf{Recall} & \textbf{Test Red.} & \textbf{Time Red.} & \textbf{Escapes} \\
\midrule
Full Suite & 100.0\% & 0.0\% & 0.0\% & 0 \\
Random-K (25\%) & 25.2\% & 75.1\% & 78.7\% & 133 \\
Changed File & 42.3\% & 83.6\% & 86.1\% & 69 \\
Static AST / Dependency & 38.6\% & 75.6\% & 68.2\% & 94 \\
Historical Failure & 48.4\% & 75.1\% & 70.0\% & 99 \\
Uncalibrated ML & 51.8\% & 75.1\% & 27.7\% & 69 \\
Calibrated No-Abstain & 51.8\% & 75.1\% & 27.7\% & 69 \\
\textbf{ConfTest (Ours)} & \textbf{100.0\%} & \textbf{3.1\%} & \textbf{0.0\%} & \textbf{0} \\
\bottomrule
\end{tabular}
\end{table}

Table~\ref{tab:baselines} reports the $183$ held-out commits. ConfTest is the only configuration with zero escaped commits, and also the one that saves least: it abstains on $97.81\%$ of commits, so its $3.1\%$ reduction in test executions yields no measurable wall-clock saving (bootstrap mean $1.05\%$, $95\%$ CI $[0.25\%, 2.10\%]$, median $-0.0\%$; $2{,}000$ resamples by commit). Every subset baseline reduces executions by $75$--$84\%$ and loses half to three quarters of the failures doing so. Cliff's $\delta$ on recall is large and positive against every baseline ($0.51$--$0.99$), and negative on time reduction against every one of them ($-0.91$ to $-1.00$); we report both directions.

Two observations deserve emphasis. First, the \emph{Uncalibrated ML} and \emph{Calibrated No-Abstain} rows are numerically identical on every selection metric. Temperature scaling is a monotone transform of the score and therefore cannot alter a top-$k$ ranking: calibration contributes nothing to selection quality and earns its place solely by making the abstention gate's confidence estimates trustworthy. Second, \emph{Uncalibrated ML} reduces test count by $75.1\%$ but wall-clock by only $27.7\%$, because the tests it ranks highest are disproportionately slow---a budget expressed in test count is not a budget in time, and RTS papers that report only the former overstate their savings.

\subsection{The Recall Floor Is Not Met at a Useful Reduction}

We pre-registered a $95\%$ held-out failure-recall floor. Selecting $\tau_{\text{abstain}}=0.044$ on the validation split yields $32.85\%$ test reduction at $96.02\%$ recall; the same threshold on the held-out split yields $32.85\%$ reduction ($95\%$ CI $[26.50, 39.18]$) at $87.69\%$ recall ($[75.77, 95.81]$), with $15$ escaped commits and $395$ missed failures. The floor is missed at the point estimate and at the CI lower bound, and the validation-to-test recall gap is $8.33$ percentage points. A post-hoc sweep of the held-out split does contain five thresholds that would clear $95\%$ (best: $26.56\%$ reduction at $97.35\%$ recall), but selecting a threshold by its test-set score would forfeit the held-out guarantee, so we record that sweep as a diagnostic and do not use it.

\subsection{Threats to Validity}

\textbf{Feature degeneracy.} Thirteen of our $32$ canonical features are constant across the training split, including all twelve code-churn columns; removing the churn group reproduces the full model exactly, and the group alone achieves ROC-AUC $0.5000$. We attribute this to a defect in diff mining rather than to a property of code churn, and no claim about churn should be drawn from our results. Static call-graph reachability is the one load-bearing family: removing it costs $15.5$ points of recall at a fixed budget.

\textbf{Stale artifacts.} Our ensemble was trained with LightGBM's \texttt{subsample} set but \texttt{subsample\_freq} left at its default of $0$, which silently disables row bagging. The seeds therefore differed only in column sampling and tie-breaking, making our epistemic $\sigma$ an underestimate of true ensemble diversity. Since $\sigma$ is precisely the quantity the abstention gate thresholds, the operating points reported here may shift once the ensemble is retrained.

\textbf{Cross-project transfer.} Leave-one-project-out over five repositories gives a macro-average recall of $51.57\%$ at a $25\%$ budget, ranging from $30.00\%$ to $88.42\%$, with ROC-AUC below chance ($0.4747$) on one project and held-out ECE reaching $0.2957$ against $0.0242$ in-distribution. Zero-shot transfer does not work; a new repository must be treated as out-of-distribution until it accumulates its own execution history.

\textbf{Latency scope.} Our reported $3.345$\,ms mean decision cost covers scoring only (sanitisation, ensemble inference, calibration, policy) and excludes diff mining, AST parsing, call-graph construction and history lookup, which dominate real per-commit cost. A cold process pays a further $1.32$\,s on its first batch.

\section{Conclusion}
ConfTest shows that epistemic uncertainty quantification and post-hoc probability calibration can support conservative machine-learned regression test selection: on $183$ held-out commits the shipped policy escaped zero failures, because it declines to select whenever ensemble disagreement or distributional novelty exceeds a tuned threshold. This is an empirical outcome, not a formal safety guarantee for future commits or unseen repositories. It also shows that safety on this sample alone does not make such a system \emph{economical}. At the threshold that holds full recall, the policy abstains on $97.81\%$ of commits and saves no measurable wall-clock time; at a threshold that saves a third of test executions, held-out recall falls to $87.69\%$ and our $95\%$ floor is violated. The binding constraint is ranker accuracy---PR-AUC $0.1393$ against a $3.71\%$ positive rate, roughly four times a no-skill baseline but far from sufficient---and not the selective-prediction machinery around it. We therefore report a negative result on the economic claim and a positive one on the observed behavior of the abstention mechanism, and we identify the gate as directly reusable by any future RTS ranker strong enough to exploit it.

\bibliographystyle{ACM-Reference-Format}
\bibliography{references}

\end{document}
